\addcontentsline{toc}{subsubsection}{Solution 23. The Euclidean--Multiplicative Squeeze}

\begin{dlsummary}[title={\bfseries Solution 23. The Euclidean--Multiplicative Squeeze},breakable]

\begin{proof}
Define
\[
g(t)=f(\sqrt t)^2\qquad(t>0).
\]
The first inequality in the statement gives
\[
g(x+y)\geq g(x)+g(y). \tag{1}
\]
For the second inequality, its right-hand side is
\[
f(xy)^2\frac{f(x)^2+f(y)^2}{f(x)^2f(y)^2}.
\]
Since $f(x)^2+f(y)^2>0$, cancellation yields
\[
f(xy)^2\leq f(x)^2f(y)^2.
\]
Applying this with $\sqrt x$ and $\sqrt y$ in place of $x$ and $y$ gives
\[
g(xy)\leq g(x)g(y). \tag{2}
\]
The composition condition is equivalently
\[
g(g(2))=f(f(\sqrt2))^2\leq2. \tag{3}
\]

\medskip
\noindent
\textbf{Step 1. Determine $g(1)$ and $g(2)$.}

Taking $x=y=1$ in (2) gives $g(1)\leq g(1)^2$. Since $g(1)>0$,
we have $g(1)\geq1$. By repeated use of (1),
\[
g(n)\geq n g(1)\geq n\qquad(n\in\mathbb N).
\]
Fix $t>1$ and a positive integer $n$. Submultiplicativity and
superadditivity imply
\[
g(t)^n\geq g(t^n)\geq\lfloor t^n\rfloor.
\]
Taking $n$th roots and letting $n\to\infty$ gives $g(t)\geq t$.
In particular, $g(2)\geq2$. Applying this bound again at $g(2)>1$,
\[
2\geq g(g(2))\geq g(2)\geq2,
\]
where (3) supplies the first inequality. Thus $g(2)=2$. Now (2) at
$x=y=1$ gives $g(1)\leq g(1)^2$, while (1) and $g(2)=2$ give
$2\geq2g(1)$. Therefore $g(1)=1$.

\medskip
\noindent
\textbf{Step 2. Obtain a global upper bound.}

For every positive integer $n$, (1) and (2) give
\[
2^n\leq g(2^n)\leq g(2)^n=2^n,
\]
so $g(2^n)=2^n$. Fix $x>0$ and, for sufficiently large $n$, let
\[
m_n=\left\lfloor\frac{2^n}{x}\right\rfloor\geq1.
\]
Write $2^n=m_nx+r_n$ with $r_n\geq0$. If $r_n>0$, repeated
superadditivity and positivity of $g(r_n)$ show that
\[
2^n=g(2^n)\geq m_ng(x)+g(r_n)\geq m_ng(x).
\]
If $r_n=0$, repeated superadditivity directly gives the same final
bound. Hence
\[
g(x)\leq\frac{2^n}{\lfloor 2^n/x\rfloor}\longrightarrow x,
\]
and so $g(x)\leq x$ for every $x>0$.

\medskip
\noindent
\textbf{Step 3. Force equality.}

By (2), $g(1)\leq g(x)g(1/x)$. Using $g(1)=1$ and the upper bound
just proved,
\[
1\leq g(x)g(1/x)\leq\frac{g(x)}x.
\]
Thus $g(x)\geq x$, and consequently $g(x)=x$ for every $x>0$.
Since $g(x)=f(\sqrt x)^2$ and $f$ is positive, it follows that
\[
f(x)=x\qquad(x>0).
\]

\medskip
\noindent
\textbf{Verification.}

For $f(x)=x$, both inequalities hold with equality, and
$f(f(\sqrt2))=\sqrt2$. Therefore this function is indeed the unique
solution.
\end{proof}

\end{dlsummary}